% ============================================================
% Section 4: Experiments and Results
% ============================================================

\section{Experiments and Results}
\label{sec:experiments}

\subsection{Experimental Setup}
\label{subsec:setup}

All experiments are conducted on a single workstation running Windows~11,
equipped with an AMD Ryzen~5 3400G (3.70\,GHz), 32\,GB DDR4-3200 RAM, and an NVIDIA
RTX~5060\,Ti (16\,GB VRAM), using PyTorch 2.10.0 with CUDA 13.0 and
Optuna~3.x. Training uses a batch size of 64 and no data augmentation
\rev{(applying standard augmentation, rotation $\pm10^\circ$, brightness,
contrast, Gaussian blur and JPEG compression, significantly degraded every
metric on the FantasyID holdout; see~\citep{docverify_zenodo})}.
All images are resized to $224\!\times\!224$ pixels and normalised to
$[0,1]$ prior to caching in GPU memory.
The full NCV pipeline required
approximately 59~hours; VRAM caching reduces per-trial training time from
$\sim$12 to $\sim$3.5 minutes.

\subsection{Nested Cross-Validation Results}
\label{subsec:ncv_results}

Table~\ref{tab:nested_cv} reports per-fold test metrics. PR-AUC is
consistently high across all folds ($0.977$--$0.999$, mean $0.9921 \pm
0.0058$). Dice shows slightly more fold-to-fold variability ($0.800$--$0.909$,
mean $0.856 \pm 0.030$), reflecting the difficulty of pixel-level localization.
The Pareto-selected $\lambda_{\text{mask}}$ clusters between 2.3 and 2.7 in
eight of ten folds; the two outlier folds (F5, F9) correspond to the two lowest
Dice scores, confirming that $\lambda_{\text{mask}}$ is a decisive driver of
localization quality.

\begin{table}[!t]
  \centering
  \caption{Per-fold test results from 10-fold NCV.
           $\overline{w}_{\text{mask}}$: selected loss weight.}
  \label{tab:nested_cv}
  \begin{tabular}{cccc}
    \toprule
    Fold & PR-AUC & Dice & $\overline{w}_{\text{mask}}$ \\
    \midrule
    F1  & 0.9939 & 0.861 & 2.46 \\
    F2  & 0.9769 & 0.871 & 2.70 \\
    F3  & 0.9993 & 0.909 & 2.47 \\
    F4  & 0.9915 & 0.858 & 2.43 \\
    F5  & 0.9948 & 0.827 & 1.41 \\
    F6  & 0.9936 & 0.874 & 2.89 \\
    F7  & 0.9940 & 0.861 & 2.67 \\
    F8  & 0.9967 & 0.887 & 2.29 \\
    F9  & 0.9896 & 0.807 & 1.95 \\
    F10 & 0.9909 & 0.800 & 2.65 \\
    \midrule
    \textbf{Mean (Std)} & \textbf{0.9921 (0.0058)} & \textbf{0.856 (0.030)} & \textbf{2.39 (0.43)} \\
    \bottomrule
  \end{tabular}
\end{table}

\subsection{Final Model Configuration}
\label{subsec:final_config}

\rev{The cross-validation above yields ten fold-specific configurations and
estimates generalization, but the deployed system is a single model. We select
the configuration with the smallest inner-CV distance to the ideal point
($1,1$) across all outer folds (Eq.~\ref{eq:ideal_point};
Table~\ref{tab:hpo_space}), which fixes both
the architecture (decoder width) and the Pareto-optimized loss weight
$\lambda_{\text{mask}}=2.46$, consistent with the $2.3$--$2.7$ cluster of
Section~\ref{subsec:ncv_results}. This final configuration is retrained from
scratch on the development set (up to 100 epochs with early stopping, as in
Section~\ref{subsec:ncv}) and evaluated on the held-out 15\% blind-test
set; we repeat the retraining over 30 independent seeds to quantify
run-to-run variability.}

\subsection{Ablation and Selection-Criterion Study}
\label{subsec:ablation}

\begin{table*}[!t]
  \centering
  \caption{\rev{Ablation and selection-criterion comparison (blind test,
           30 seeds); all variants vs.\ \textbf{Pareto (ours)}. $p$: paired
           Wilcoxon~\citep{wilcoxon1945}\,+\,Holm~\citep{holm1979} (significant
           at $p<0.05$); $d$: Cohen's~$d$~\citep{cohen1988statistical}
           ($|d|\!\approx\!0.2$ small, $0.5$ medium, $\geq\!0.8$ large).}}
  \label{tab:ablation}
  \begin{tabular}{lcccccc}
    \toprule
    & \multicolumn{3}{c}{PR-AUC} & \multicolumn{3}{c}{Dice} \\
    \cmidrule(r){2-4}\cmidrule(l){5-7}
    Variant & Mean $\pm$ Std & $p$ & $d$ & Mean $\pm$ Std & $p$ & $d$ \\
    \midrule
    \textbf{Pareto (ours)} & $\mathbf{0.9967 \pm 0.0021}$ & — & — & $\mathbf{0.875 \pm 0.018}$ & — & — \\
    \rev{\texttt{by\_prauc}} & \rev{$0.9953 \pm 0.0021$} & \rev{$0.06$} & \rev{$0.49$} & \rev{$0.859 \pm 0.023$} & \rev{$0.04$} & \rev{$0.59$} \\
    \rev{\texttt{by\_dice}}  & \rev{$0.9948 \pm 0.0027$} & \rev{$0.007$} & \rev{$0.64$} & \rev{$0.857 \pm 0.015$} & \rev{$<0.001$} & \rev{$1.00$} \\
    \midrule
    \texttt{cls\_only}          & $0.9503 \pm 0.0744$ & $<0.001$ & $0.62$ & $0.002 \pm 0.008$ & $<0.001$ & $47.2$ \\
    \texttt{seg\_only}          & $0.7277 \pm 0.0528$ & $<0.001$ & $5.15$ & $0.867 \pm 0.023$ & $0.29$ & $0.32$ \\
    \texttt{unweighted\_losses} & $0.9958 \pm 0.0016$ & $0.08$ & $0.37$ & $0.857 \pm 0.016$ & $<0.001$ & $1.01$ \\
    \bottomrule
  \end{tabular}
\end{table*}

To isolate each design decision, we evaluate this final model against five
variants under the same 30-seed blind-test protocol
(Table~\ref{tab:ablation}). Three variants probe the architecture and loss
weighting: \texttt{cls\_only} (classification head only); \texttt{seg\_only}
(segmentation head only, detection by thresholding mean mask probability);
and \texttt{unweighted\_losses} (full multitask, $\lambda_{\text{cls}} =
\lambda_{\text{mask}} = 0.5$, bypassing HPO). \rev{Two further variants probe
the HPO selection criterion applied to the \emph{same} 500 MOTPE trials:
\texttt{by\_prauc} (the trial maximizing validation PR-AUC) and
\texttt{by\_dice} (the trial maximizing validation Dice), against our Pareto
criterion (distance to the ideal point).}

Regarding the architecture and loss variants, three findings stand out.
First, joint multitask training is essential:
\texttt{cls\_only} produces near-zero Dice ($0.002 \pm 0.008$). Second, a
dedicated classification head is necessary: \texttt{seg\_only} achieves
competitive Dice ($0.867$) but PR-AUC drops to $0.728 \pm 0.053$, confirming
that the segmentation logit alone is an unreliable detection signal. Third,
Pareto-optimized weights outperform equal weights specifically on localization:
\texttt{unweighted\_losses} matches the full model on PR-AUC ($p = 0.083$) but
shows significantly lower Dice ($p = 1.3 \times 10^{-4}$, $d = 1.01$, large
effect), validating the HPO contribution.

A key question for any multi-objective HPO framework is whether the added
complexity over simpler scalar alternatives yields measurable gains. \rev{Pareto
significantly outperforms both scalar criteria after Holm correction:
it improves localization over \texttt{by\_prauc} (Dice $+0.016$, $d=0.59$) and
improves \emph{both} objectives over \texttt{by\_dice} (PR-AUC $+0.0019$,
$d=0.64$; Dice $+0.018$, $d=1.00$, large effect).}

Beyond predictive accuracy, Pareto selection requires no a priori commitment to a
primary objective, making the full trade-off landscape observable
(Fig.~\ref{fig:pareto}) and remaining structurally consistent with the DeepID
challenge protocol, which ranks detection and localization independently.

\subsection{Comparison with Related Systems}
\label{subsec:challenge}

\rev{DocVerify was developed after the DeepID 2025
Challenge~\citep{deepid2025challenge} deadline, so official test-set
evaluation was not possible; we report results on our internal 15\% holdout
(30 seeds), which differs from the official challenge set (FantasyID test
partition plus a private PXL~Vision set), making any comparison indicative
rather than controlled. Training data also differ markedly. The leading
entries rely on external data: Sunlight (1st) augmented FantasyID with over
60\,000 identity images, while the third-place EdgeDoc+TruFor
fusion~\citep{edgedoc2025} and the TruFor baseline~\citep{trufor2023} build on
a backbone pre-trained on $\sim$828\,000 manipulation images. EdgeDoc
(standalone), UAM-Biometrics, and DocVerify instead train on FantasyID alone,
the most directly comparable group. With these caveats, DocVerify reaches
detection F1@0.5 \rev{(the F1 score of the bona-fide/attack decision at
classification threshold $0.5$)} $=0.969 \pm 0.014$ and localization
F1$_{\text{loc}}$ \rev{(pixel-level F1 of the predicted mask)}
$= 0.807 \pm 0.096$, in line with the challenge winner
Sunlight (F1@0.5 $0.991$, F1$_{\text{loc}}$ $0.784$) and ahead of the
FantasyID-only entries (EdgeDoc $0.43$, UAM-Biometrics $0.712$), despite using
no external data; we make no superiority claim given the differing evaluation
sets.}

\subsection{\rev{Generalization}}
\label{subsec:generalization}

\rev{Two experiments probe whether the approach generalizes beyond its specific
setting. \emph{Across architectures}, applying the identical MOTPE/Pareto protocol
to an ImageNet-pretrained ResNet-18 backbone yields comparable blind-test
performance (PR-AUC $0.994$, Dice $0.889$; archived
externally~\citep{docverify_zenodo}), with neither model dominating, so the
approach generalizes across architectures. \emph{Across datasets}, evaluating the blind-test models zero-shot on the
out-of-domain SIDTD dataset~\citep{sidtd2024} (templates subset; 1\,000 bonafide,
1\,222 attack) drops performance to near chance (PR-AUC $0.555 \pm 0.011$, ROC-AUC
$0.504 \pm 0.009$): the FantasyID$\rightarrow$SIDTD domain gap is too wide for a
model trained on FantasyID alone. Retraining the same architecture on the SIDTD
training partition, however, recovers strong performance (PR-AUC $0.878 \pm 0.016$,
ROC-AUC $0.826 \pm 0.023$), showing the gap reflects domain shift rather than
architectural capacity; the zero-shot scope of the current FantasyID-trained system
is thus bounded to its training domain, motivating multi-source training as future
work.}

\subsection{Qualitative Analysis}
\label{subsec:qualitative}

Figure~\ref{fig:qualitative} shows a representative output for a forged
document with five altered regions. The predicted mask successfully localizes
all five tampered fields — photograph, given name, date of birth, document
number, and expiry date — with tight spatial alignment to the ground-truth
alterations.

\begin{figure*}[!t]
  \centering
  \begin{tabular}{@{}ccc@{}}
    \small\textbf{Input document} &
    \small\textbf{Predicted mask} &
    \small\textbf{Overlay} \\[2pt]
    \includegraphics[width=0.318\textwidth]{fig_qual_attack_input} &
    \includegraphics[width=0.318\textwidth]{fig_qual_attack_mask}   &
    \includegraphics[width=0.318\textwidth]{fig_qual_attack_overlay} \\[1pt]
    \multicolumn{3}{c}{\small Forged document — five altered regions localized} \\
  \end{tabular}
  \caption{Qualitative result on a forged FantasyID sample. \textbf{Left:}
  input document. \textbf{Centre:} binary mask predicted by DocVerify.
  \textbf{Right:} predicted mask overlaid in red on the input. All five
  tampered fields are correctly localized.}
  \label{fig:qualitative}
\end{figure*}

